% 来源：成文/A\_06\_模型假设.md
\section{模型假设}
本文按几何与介质、物性与耦合、表面与环境、几何演化与适用域四组给出八条假设。每条给出合理性依据（为什么可以这样假设）与适用前提（什么情况下不成立），后者同时说明假设失效时的后果。

第一组为几何与介质前提。

\textbf{假设 1}　药材视为无限长圆柱，温度场与干基含水率场只沿半径方向与时间变化。依据：长径比为 6.25，两端面面积仅占圆柱总表面积约 7.4\%，端面贡献小\textsuperscript{\cite{ref14}}；此处须一并声明所引判据的适用落差——文献\textsuperscript{\cite{ref14}}对圆柱试样给出的长径比判据为 $L_0/R_0>15$（此时一维近似对扩散系数的高估误差小于 5\%），本题 $L_0/R_0$ 为 12.5，\textbf{略低于}该判据，故端面效应按已声明的方向处理，其误差方向与量级见第七章不足 2。适用前提：轴向近似均匀；端面效应显著时须改用二维模型，否则轴向中部的结果将偏低。

\textbf{假设 2}　药材均质且各向同性，密度、比热容、导热系数与扩散系数均取标量。依据：题面按标量给出四个物性经验式，未给任何方向相关参数。适用前提：视为连续均质介质；强各向异性物料不适用，须改写为张量物性形式。

第二组为物性与耦合口径。

\textbf{假设 3}　物性经验式按问题分套采用：问题一用附录 2，问题二、三全程用附录 3，问题四用附录 4；各套经验式在本题实际使用范围内视为有效。依据：题面明确规定各问所用附录，并说明经验式统一采用附录 3。适用前提：干基含水率由 2.55 降至 0.15、温度约 28 至 50 ℃；低含水率端的外推代价已在第六章量化\textsuperscript{\cite{ref5,ref6}}，若经验式在该端失真，尾段时长将出现同量级偏差。

\textbf{假设 4}　忽略蒸发潜热与热质交叉扩散，能量方程不含源项；固相与液相处于局部热力学平衡，用单一温度描述。依据：题面未给汽化潜热、交叉扩散系数与吸湿等温线，引入将带来不可标定的参数\textsuperscript{\cite{ref3}}。适用前提：本文已对潜热作对照计算，代价量化为温度系统性偏低约 0.07 至 0.10 ℃，属已知且有界；若计入潜热，脱湿变慢、烘干时长变长。

第三组为表面与环境。

\textbf{假设 5}　表面换热系数与对流传质系数全程沿用附录 2 取值，即 25 W/(m\textsuperscript{2}·K) 与 $8\times10^{-7}$ m/s。依据：题面仅在该附录给出这两个系数，其余各问未另给取值。适用前提：恒温干燥段真实系数若不同会产生偏差，本文已在灵敏度分析中对两者作 ±20\% 扰动。

\textbf{假设 6}　表面传质以干基浓度差为驱动力，表面水分通量正比于表面含水率与烘房水分浓度之差，不引入水分活度与平衡含水率。依据：附录 2 直接以该形式给出表面边界式，且未给吸湿等温线。适用前提：低含水率段的真实驱动力若为活度差，表面通量会出现偏差，已在第七章列为局限。

\textbf{假设 7}　烘房环境单向给定：环境温度与水分浓度由附件 1 的实测数据决定，不因药材吸热放湿而改变；14400 s 之后取恒定值 50.00 ℃ 与 0.04999 kg/kg，该两值由附件 1 末段实测均值外推。依据：附件 1 为固定时序数据，题面未给药材对环境的反馈关系，且烘房相对单根药材可视为大空间。适用前提：环境作为给定边界；若需考核双向耦合，应另行联立求解。

第四组为几何演化与适用域。

\textbf{假设 8}　问题四的半径由附件 2 节点线性插值给出，259200 s 之后取常值 1.198 cm；并采用仿射（均匀）收缩，即物质点半径按比例缩放，以 $\xi$ 为物质坐标。依据：附件 2 只给表面半径时序，题面未规定内部收缩分配方式；仿射收缩是无需额外参数的最简分配，等价于各向同性率为 1 的特例，与整体失水导致均匀收缩的现象相符\textsuperscript{\cite{ref4,ref7}}。适用前提：忽略表层先干先缩的非均匀效应；若真实收缩非均匀，问题四的内部水分场与烘干时长将出现偏差，本文未做该项对照并已列为改进方向。

说明：题面笔误与数据换算（如附录 3、附录 4 中扩散系数经验式的温度单位）属更正题目，不计入本文假设，正文在相应位置随文说明。
